\documentclass[10pt,a4paper]{amsart}
\usepackage[margin=19mm]{geometry}
\usepackage{amsmath,amssymb,mathtools}
\usepackage[scr=rsfs,cal=euler]{mathalfa}
\usepackage{xfrac}
\usepackage[unicode,hidelinks,bookmarks=false]{hyperref}
\newcommand{\ie}{i.e.\ }
\newcommand{\cf}{cf.\ }
\newcommand{\resp}{respectively\ }
\newcommand{\Subsection}{\subsection}
\providecommand{\nobreakdash}{\nobreak\hbox{-}}
\setlength{\emergencystretch}{2em}
\multlinegap0pt
\usepackage{amssymb}
\usepackage{mathtools}
\usepackage{bm}
\usepackage{algorithm}\usepackage{algpseudocode}
\usepackage{tcolorbox}
\newtheorem{assump}{Assumption}
\newtheorem{lemma}{Lemma}
\newcommand{\dd}{\,\mathrm{d}}
\title{Mirror Descent for Deterministic Optimal Control}
\author{Ye Feng, Jianfeng Lu}
\date{Source: arXiv:2605.02653v2 (CC BY 4.0)}
\begin{document}
\maketitle

\noindent\emph{Source and licence.} Mirror Descent for Deterministic Optimal Control. Version arXiv:2605.02653v2 (CC BY 4.0), distributed by arXiv under the Creative Commons Attribution 4.0 International licence, http://creativecommons.org/licenses/by/4.0/, as recorded in the arXiv licence metadata for this exact version and shown on the abstract page https://arxiv.org/abs/2605.02653v2. Full licence notice: \texttt{licenses/arxiv-2605.02653-CC-BY-4.0.txt}.\par\smallskip

\noindent\textit{Editorial scope.} Counted unit: the article's lemma lem:first-variation (source0.tex lines 240--262). The article's complete original proof of it is delivered with it (source0.tex lines 1630--1987, one \texttt{proof} environment of 358 lines, which the article places away from the statement). No mathematics is written, repaired or re-derived: the statement and the proof are the article's own lines, in the article's own notation, and no retained line is substituted or re-flowed.  Retained uncounted as the source-local context the proof needs: lines 121--150.  Nothing was omitted from any retained span, so the package carries no omission record.  No retained line contains a citation, so the wrapper carries no bibliography and no bibliography entry.  No retained line contains a figure environment, an image-inclusion command or drawing code, so no image is delivered and the package declares no asset dependency.  Editorial formatting only, in addition: an AMS \texttt{amsart} preamble that restates the article's own declarations and macros used by the retained lines, namely the environments \texttt{assump}, \texttt{lemma} on separate counters, together with the standard packages those lines need.  The retained environments print in this document as Assump 1, Lemma 1 in order of appearance, whereas the article prints the same items with the numbering of its own full document.  The complete native source of the article as distributed in the arXiv e-print for this exact version is bundled unchanged as \texttt{sources/199733/source0.tex}.
\begin{assump}\label{asp:main_asp}
    We make the following standing assumptions:
    \begin{enumerate}
        \item $U$ is a nonempty, closed, convex subset of $\mathbb{R}^m$.
        \item The maps $b: [0,T] \times \mathbb{R}^d \times U \to \mathbb{R}^d$ and $f: [0,T] \times \mathbb{R}^d \times U\to \mathbb{R}$ are continuous, and there exists a constant $M > 0$ such that for $\varphi = b, f$,
              \begin{align}
                  |\varphi_t(x, u) - \varphi_t(x', u')| \leq M(|x - x'| + |u - u'|), \quad |\varphi_t(0, u)| \leq M, \nonumber
              \end{align}
              for any $t \in [0,T]$, $x, x' \in \mathbb{R}^d$, and $u, u' \in U$.
        \item For $\varphi = b, f$, the map $\varphi_t(x,u)$ is twice continuously differentiable in $(x,u)$, and, possibly after increasing $M$, the same constant satisfies
              \begin{align}
                  |\nabla_x \varphi_t(x, u) - \nabla_x \varphi_t(x', u')| \leq M(|x - x'| + |u - u'|), \nonumber
              \end{align}
              for any $t \in [0,T]$, $x, x' \in \mathbb{R}^d$, and $u, u' \in U$.
              Further, assume that there exist constants $M_{b,uu}, M_{f,uu} > 0$ such that
\[
\|\nabla^2_{uu} b_t(x,u)\| \le M_{b,uu},
\qquad
\|\nabla^2_{uu} f_t(x,u)\| \le M_{f,uu},
\]
for all $(t,x,u) \in [0,T]\times\mathbb{R}^d\times U$.
        \item The map $g: \mathbb{R}^d \to \mathbb{R}$ is continuously differentiable and satisfies
        \[
        \max\{|g(x) - g(x')|, |\nabla g(x) - \nabla g(x')|\} \leq M|x - x'|, \quad |g(0)| \leq M,
        \]
        for any $x, x' \in \mathbb{R}^d$.
        \item The regularizer $h : U \to \mathbb{R}$ is convex and continuously differentiable.
        \item There exists a (continuous) optimal solution to Problem (\textbf{P}).
    \end{enumerate}
\end{assump}

\begin{lemma}[First variation formula]\label{lem:first-variation}
Under Assumption~\ref{asp:main_asp}, for every $u,v\in\mathcal U$,
the one-sided directional derivative of $J^0$ at $u$ along the feasible segment from $u$ to $v$ exists and is given by
\begin{equation}\label{eq:first-var-J0}
dJ^0(u)(v-u)
=
-\int_0^T \nabla_u \mathcal H_t^0(x_t,p_t,u_t)\cdot (v_t-u_t)\dd t,
\end{equation}
where $(x, p)$ is the state-adjoint pair corresponding to $u$.

Moreover, the corresponding one-sided directional derivative of $J^\tau$ is
\begin{equation}\label{eq:first-var-Jtau}
dJ^\tau(u)(v-u)
=
- \int_0^T \nabla_u \mathcal H_t^\tau(x_t,p_t,u_t)\cdot (v_t-u_t)\dd t,
\end{equation}
where
\[
\nabla_u \mathcal H_t^\tau(x,p,u)
=
\nabla_u \mathcal H_t^0(x,p,u) - \tau \nabla h(u).
\]
\end{lemma}

\begin{proof}[Proof of Lemma~\ref{lem:first-variation}]
Fix $u,v\in\mathcal U$ and set
\[
\delta u := v-u.
\]
For $\varepsilon \in [0,1]$, define the perturbed control
\[
u^\varepsilon := u+\varepsilon \delta u.
\]
Since $U$ is convex, $u^\varepsilon_t\in U$ for every $t$ whenever $\varepsilon\in[0,1]$,
hence $u^\varepsilon\in\mathcal U$. We compute the one-sided directional derivative at $\varepsilon=0$ along this feasible segment.

Let $x$ and $x^\varepsilon$ denote the states associated with $u$ and $u^\varepsilon$ respectively:
\begin{align*}
& \dot x_t=b_t(x_t,u_t),
\qquad
x(0)=x_0; \\
& \dot x_t^\varepsilon=b_t(x_t^\varepsilon,u_t^\varepsilon),
\qquad
x^\varepsilon(0)=x_0.
\end{align*}

We first identify the first-order variation of the state.
Define $y$ as the solution of the linearized equation
\begin{equation}\label{eq:linearized-state}
\dot y_t
=
\nabla_x b_t(x_t,u_t)\,y_t+\nabla_u b_t(x_t,u_t)\,\delta u_t,
\qquad
y_0=0.
\end{equation}

We claim that
\begin{equation}\label{eq:state-diff-quotient}
\frac{x^\varepsilon-x}{\varepsilon}\to y
\qquad\text{in } C([0,T];\mathbb R^d)
\quad\text{as }\varepsilon\to 0^+.
\end{equation}

To prove this, define the remainder
\[
r_t^\varepsilon
:=
x_t^\varepsilon-x_t-\varepsilon y_t.
\]
Then $r_0^\varepsilon=0$, and
\[
\dot r_t^\varepsilon
=
b_t(x_t^\varepsilon,u_t^\varepsilon)-b_t(x_t,u_t)
-
\varepsilon\Big(
\nabla_x b_t(x_t,u_t)y_t+\nabla_u b_t(x_t,u_t)\delta u_t
\Big).
\]
Now write
\[
x_t^\varepsilon=x_t+\varepsilon y_t+r_t^\varepsilon,
\qquad
u_t^\varepsilon=u_t+\varepsilon\delta u_t.
\]
Using the first-order Taylor expansion of $b_t(\cdot,\cdot)$ in the variables $(x,u)$
around $(x_t,u_t)$, we obtain
\[
b_t(x_t^\varepsilon,u_t^\varepsilon)-b_t(x_t,u_t)
=
\nabla_x b_t(x_t,u_t)\bigl(\varepsilon y_t+r_t^\varepsilon\bigr)
+
\nabla_u b_t(x_t,u_t)\,\varepsilon\delta u_t
+
\rho_t^\varepsilon,
\]
where the remainder $\rho_t^\varepsilon$ satisfies
\[
|\rho_t^\varepsilon|
\le
C\Big(|\varepsilon y_t+r_t^\varepsilon|+|\varepsilon\delta u_t|\Big)^2
\]
for a constant $C>0$ independent of $\varepsilon$ and $t$.
Hence
\[
\dot r_t^\varepsilon
=
\nabla_x b_t(x_t,u_t) r_t^\varepsilon + \rho_t^\varepsilon.
\]
Therefore
\begin{equation}\label{eq:re_e}
    |r_t^\varepsilon|
\le
\int_0^t |\nabla_x b_s(x_s,u_s)|\,|r_s^\varepsilon|\dd s
+
\int_0^t |\rho_s^\varepsilon|\dd s.
\end{equation}
By Assumption~\ref{asp:main_asp}, $\sup_{s \in [0,T]} |\nabla_x b(x_s, u_s)|\leq M$.
Moreover,
\begin{equation}\label{eq:rho_e}
    |\rho_t^\varepsilon|
\le
C\Bigl(\varepsilon^2|y_t|^2+|r_t^\varepsilon|^2+\varepsilon^2|\delta u_t|^2\Bigr).
\end{equation}
We estimate the square term $|r_t^{\varepsilon}|^2$ as follows. By the definition of $r_t^{\varepsilon}$,
\[
\sup_{t\in[0,T]} |r_t^\varepsilon|
\le
\sup_{t\in[0,T]} |x_t^\varepsilon-x_t|
+
\varepsilon\,\sup_{t\in[0,T]} |y_t|.
\]
By Assumption~\ref{asp:main_asp} (2) and $x_0^{\varepsilon} = x_0$,
\[
|x_t^\varepsilon - x_t| \leq \int_0^t |b_s(x_s^\varepsilon, u_s^\varepsilon) - b_s(x_s, u_s)| \dd s \leq \int_0^t M \left( |x_s^\varepsilon - x_s| + |u_s^\varepsilon - u_s| \right) \dd s.
\]
By Gr\"onwall's inequality,
\[
\sup_{t\in[0,T]} |x_t^\varepsilon-x_t|
\le
M e^{MT} \|u^\varepsilon-u\|_{L^1}
=
M e^{MT} \varepsilon\,\|\delta u\|_{L^1},
\]
By the linearized equation \eqref{eq:linearized-state} and $y_0 = 0$,
\[
|y_t| \leq \int_0^t |\nabla_x b_s(x_s, u_s)|\,|y_s|\dd s + \int_0^t |\nabla_u b_s(x_s, u_s)| \, |\delta u_s| \dd s.
\]
Assumption~\ref{asp:main_asp} yields
\[
|y_t| \leq M \left( \int_0^t |y_s| \dd s + \int_0^t |\delta u_s| \dd s \right).
\]
By Gr\"onwall's inequality,
\[
\sup_{t \in [0,T]} |y_t| \leq M e^{MT}\|\delta u\|_{L^1}.
\]
Therefore,
\[
\sup_{t\in[0,T]} |r_t^\varepsilon|
\le
C' \varepsilon, \qquad C':= 2 Me^{MT} \|\delta u\|_{L^1}.
\]
Substituting back to \eqref{eq:rho_e} and
using that $y\in C([0,T];\mathbb R^d)$ and $\delta u\in L^2([0,T];\mathbb R^m)$, we obtain
\[
\int_0^T |\rho_t^\varepsilon| \dd t = o(\varepsilon).
\]
By applying Gr\"onwall to \eqref{eq:re_e},
\[
\sup_{t\in[0,T]} \left|\frac{r_t^\varepsilon}{\varepsilon}\right|\to 0,
\]
which proves \eqref{eq:state-diff-quotient}.

We next compute the derivative of the cost.
By definition,
\[
J^0(u^\varepsilon)-J^0(u)
=
\int_0^T \Big(f_t(x_t^\varepsilon,u_t^\varepsilon)-f_t(x_t,u_t)\Big)\dd t
+
g(x_T^\varepsilon)-g(x_T).
\]
We treat the running cost and terminal cost separately.

For the running cost, by the first-order Taylor expansion of $f_t(\cdot,\cdot)$
around $(x_t,u_t)$,
\[
f_t(x_t^\varepsilon,u_t^\varepsilon)-f_t(x_t,u_t)
=
\nabla_x f_t(x_t,u_t)\cdot (x_t^\varepsilon-x_t)
+
\nabla_u f_t(x_t,u_t)\cdot (u_t^\varepsilon-u_t)
+
\tilde\rho_t^\varepsilon,
\]
where
\[
|\tilde\rho_t^\varepsilon|
\le
C\Big(|x_t^\varepsilon-x_t|+|u_t^\varepsilon-u_t|\Big)^2.
\]
Since
\[
x_t^\varepsilon-x_t=\varepsilon y_t+r_t^\varepsilon,
\qquad
u_t^\varepsilon-u_t=\varepsilon \delta u_t,
\]
and
\[
\sup_{t\in[0,T]}|r_t^\varepsilon|=o(\varepsilon),
\]
it follows that
\[
\int_0^T |\tilde\rho_t^\varepsilon|\dd t=o(\varepsilon).
\]
Hence
\[
\int_0^T \Big(f_t(x_t^\varepsilon,u_t^\varepsilon)-f_t(x_t,u_t)\Big)\dd t
=
\varepsilon\int_0^T
\Big(
\nabla_x f_t(x_t,u_t)\cdot y_t
+
\nabla_u f_t(x_t,u_t)\cdot \delta u_t
\Big)\dd t
+o(\varepsilon).
\]

For the terminal cost, since $g$ is continuously differentiable,
\[
g(x_T^\varepsilon)-g(x_T)
=
\nabla g(x_T)\cdot (x_T^\varepsilon-x_T)+o(|x_T^\varepsilon-x_T|).
\]
Using
\[
x_T^\varepsilon-x_T=\varepsilon y_T+r_T^\varepsilon
\qquad\text{and}\qquad
r_T^\varepsilon=o(\varepsilon),
\]
we obtain
\[
g(x_T^\varepsilon)-g(x_T)
=
\varepsilon\,\nabla g(x_T)\cdot y_T+o(\varepsilon).
\]

Combining the two expansions, we conclude that
\[
J^0(u^\varepsilon)-J^0(u)
=
\varepsilon\int_0^T
\Big(
\nabla_x f_t(x_t,u_t)\cdot y_t
+
\nabla_u f_t(x_t,u_t)\cdot \delta u_t
\Big)\dd t
+
\varepsilon\,\nabla g(x_T)\cdot y_T
+
o(\varepsilon).
\]
Therefore,
\begin{equation}\label{eq:pre-adjoint-first-variation}
dJ^0(u)(\delta u)
=
\int_0^T
\Big(
\nabla_x f_t(x_t,u_t)\cdot y_t
+
\nabla_u f_t(x_t,u_t)\cdot \delta u_t
\Big)\dd t
+
\nabla g(x_T)\cdot y_T.
\end{equation}

It remains to eliminate the state variation $y$ using the adjoint equation.
Recall that the adjoint $p$ satisfies
\[
\dot p_t
=
-\nabla_x \mathcal H_t^0(x_t,p_t,u_t),
\qquad
p_T= - \nabla g(x_T).
\]
Since
\[
\mathcal H_t^0(x,p,u)=p\cdot b_t(x,u) - f_t(x,u),
\]
we have
\[
\dot p_t
=
-\nabla_x b_t(x_t,u_t)^\top p_t + \nabla_x f_t(x_t,u_t).
\]
We now compute
\[
\frac{d}{dt}(p_t\cdot y_t)=\dot p_t\cdot y_t+p_t\cdot \dot y_t.
\]
Using the equations for $p$ and $y$, we get
\begin{align*}
\frac{d}{dt}(p_t\cdot y_t)
&=
\Big(-\nabla_x b_t(x_t,u_t)^\top p_t + \nabla_x f_t(x_t,u_t)\Big)\cdot y_t \\
&\quad
+
p_t\cdot
\Big(
\nabla_x b_t(x_t,u_t)y_t+\nabla_u b_t(x_t,u_t)\delta u_t
\Big).
\end{align*}
The two terms involving $\nabla_x b_t$ cancel, and thus
\[
\frac{d}{dt}(p_t\cdot y_t)
=
\nabla_x f_t(x_t,u_t)\cdot y_t
+
p_t\cdot \nabla_u b_t(x_t,u_t)\delta u_t.
\]
Integrating over $[0,T]$ and using $y_0=0$ and $p_T= - \nabla g(x_T)$, we obtain
\[
- \nabla g(x_T)\cdot y_T
=
\int_0^T \nabla_x f_t(x_t,u_t)\cdot y_t\dd t
+
\int_0^T p_t\cdot \nabla_u b_t(x_t,u_t)\delta u_t\dd t.
\]
Substituting this identity into \eqref{eq:pre-adjoint-first-variation}, the
terms involving $\nabla_x f_t(x_t,u_t)\cdot y_t$ cancel and we obtain
\[
dJ^0(u)(\delta u)
=
\int_0^T
\Big(
- p_t\cdot \nabla_u b_t(x_t,u_t)\delta u_t
+
\nabla_u f_t(x_t,u_t)\cdot \delta u_t
\Big)\dd t.
\]
Equivalently,
\[
dJ^0(u)(\delta u)
=
\int_0^T
\Big(
-\nabla_u b_t(x_t,u_t)^\top p_t+\nabla_u f_t(x_t,u_t)
\Big)\cdot \delta u_t\dd t.
\]
Since
\[
\nabla_u \mathcal H_t^0(x_t,p_t,u_t)
=
\nabla_u b_t(x_t,u_t)^\top p_t - \nabla_u f_t(x_t,u_t),
\]
we conclude that
\[
dJ^0(u)(\delta u)
=
- \int_0^T \nabla_u \mathcal H_t^0(x_t,p_t,u_t)\cdot \delta u_t\dd t.
\]
Recalling that $\delta u=v-u$, this proves \eqref{eq:first-var-J0}.

Finally, since
\[
J^\tau(u)=J^0(u)+\tau \int_0^T h(u_t)\dd t,
\]
and $h$ is continuously differentiable, we have
\[
dJ^\tau(u)(\delta u)
=
dJ^0(u)(\delta u)
+
\tau \int_0^T \nabla h(u_t)\cdot \delta u_t\dd t.
\]
Therefore
\[
dJ^\tau(u)(\delta u)
=
- \int_0^T \nabla_u \mathcal H_t^\tau(x_t,p_t,u_t)\cdot \delta u_t\dd t,
\]
which proves \eqref{eq:first-var-Jtau}.
\end{proof}

\end{document}
